\documentclass[11pt]{article}
\usepackage{fontspec}
\usepackage[french,english]{babel}
\usepackage[leqno]{amsmath}
\usepackage{amssymb}
\usepackage{amsthm}
\usepackage{mathrsfs}
\usepackage{tikz-cd}
\usepackage{enumitem}

\newcommand{\printedpage}[1]{\marginpar{\footnotesize p.~#1}}
\newcommand{\msp}[1]{\marginpar{\footnotesize [#1]}}

\newcommand{\cE}{\mathcal{E}}
\newcommand{\cC}{\mathcal{C}}
\newcommand{\cM}{\mathcal{M}}
\newcommand{\Ob}{\operatorname{Ob}}
\newcommand{\Fl}{\operatorname{Fl}}
\newcommand{\Hom}{\operatorname{Hom}}
\newcommand{\End}{\operatorname{End}}
\newcommand{\Ind}{\operatorname{Ind}}
\newcommand{\Ens}{\mathrm{Ens}}
\newcommand{\pr}{\mathrm{pr}}
\newcommand{\id}{\mathrm{id}}
\newcommand{\card}{\operatorname{card}}
\newcommand{\cardred}{\operatorname{card}\,\text{\upshape red}}
\newcommand{\defeq}{\overset{\text{\upshape def}}{=}}

\begin{document}

\setcounter{footnote}{174}

%=====================================================================
\printedpage{159}
% head: A. Grothendieck : Dérivateurs, Ch. XVIII, Édition des Universités de Montpellier II et Paris VII

\msp{page 419}
is $\pi$-accessible, when $x_0$ is in $\cE_{0\,\pi}$. Now by the hypothesis on $\pi$, it is $\pi$-accessible if $x_0$ is in $\cE_{0\,\pi_0}$. But if $x$ is in $\cE_{0\,\pi}$, we will have % [sic] x -> x_0, the formula that follows concerns x_0
\[
x_0 \;=\; \varinjlim_{i\in I} x_{0\,i} \qquad
\left\{
\begin{array}{l}
I\ \ \text{large}\ \ \text{relative to}\ \ \pi_0,\\
\text{the}\ \ x_{0\,i}\ \ \text{in}\ \ \cE_{0\,\pi_0},\\
\card I \leqslant \pi
\end{array}
\right.
\]
(for $\pi_0$ is useful for $\cE_0$ by hypothesis, and $\pi$ satisfies $\pi^{\pi_0} = \pi \geqslant \cardred \cE_{\pi_0}$). As $x_0 \longrightarrow H(x_0,x_1) : \cE_0 \longrightarrow \Ens$ satisfies $\mathrm{L}_{\pi_0}$, for every $x_1$ in $\cE_1$, we will therefore have % [sic] x_0 \longrightarrow H(x_0,x_1) -> x_0 \longmapsto H(x_0,x_1), elsewhere such functors are written with \longmapsto
\[
H(x_0) \;\simeq\; \varprojlim_{i\in I} H(x_{0\,i})
\]
(where for $x_0$ in $\cE_0$, $H(x)$ denotes the functor $x_1 \longmapsto H(x_0,x_1) : \cE_1 \longrightarrow (\Ens)$). As % [sic] H(x) -> H(x_0), the functor depends on x_0
the functors $H(x_{0\,i})$ are $\pi$-accessible, and $\card I \leqslant \pi$, it follows that their $\varprojlim\limits_{I}$ $H(x_0)$ is $\pi$-accessible (SGA 4 I 9.6 (ii)). This therefore completes the proof of 2°).

3°). To prove that $\Ind_\pi(\cE(\pi)) \longrightarrow \cE$ is essentially surjective, it suffices to see that every $x \in \Ob\cE_0$ and every $x \in \Ob\cE_1$
\msp{page 420}
is in the essential image. Now this follows from the hypothesis that $\pi$ is useful for $\cE_0$ and $\cE_1$, and from 1°) which implies that $\cE_0 \longrightarrow \cE$ and $\cE_1 \longrightarrow \cE$ satisfy $\mathrm{L}_\pi$.

This completes the proof of proposition 4.14.24, and likewise the following result, more precise than the implication c) $\Longrightarrow$ a).

\textbf{Corollary 4.14.25.} \textit{Let $\cE$ be a category, $\cE_0$ an open subcategory, $\cE_1$ the complementary closed subcategory,}
\[
H \;=\; ((x_0,x_1) \longmapsto \Hom_{\cE}(x_0,x_1)) \;:\; \cE_0^\circ \times \cE_1 \;\longrightarrow\; \Ens
\]
\textit{the gluing functor.}
\begin{enumerate}
\item[\textit{1°)}] \textit{Let $\pi$ be an infinite cardinal. In order that $\cE$ and $\cE_0$ satisfy $L_\pi$, and the inclusion functor $\cE_0 \hookrightarrow \cE$ likewise, it is necessary and sufficient that $\cE_0$ satisfy $L_\pi$, and that for every $x_1$ in $\cE_1$, the functor $x_0 \longmapsto H(x_0,x_1) : \cE_0 \longrightarrow \Ens^\circ$ satisfy $L_\pi$. Then $\cE_1$ and $\cE_1 \hookrightarrow \cE$ also satisfy $L_\pi$,}
\msp{page 421}
\textit{and we have}
\[
\cE_\pi \;\subset\; \cE(\pi) \;\defeq\;
\begin{array}{l}
\text{\itshape strictly\ \ full\ \ subcategory\ \ of\ \ $\cE$\ \ such\ \ that}\\
\Ob\cE(\pi) = \Ob\cE_{0\,\pi} \cup \Ob\cE_{1\,\pi}.
\end{array}
\]
\item[\textit{2°)}] \textit{Suppose that condition c) of 4.14.24 is satisfied, and let $\pi_0$ be an infinite cardinal such that $\pi_0$ is useful for $\cE_0$, and that the functors $x_0 \longmapsto H(x_0,x_1) : \cE_0 \longrightarrow \Ens^\circ$ defined by the $x_1 \in \Ob\cE_1$ are $\pi_0$-accessible. Let $\pi$ be a cardinal such that $\pi^{\pi_0} = \pi \geqslant \sup(\pi_0, \cardred\cE_{0\,\pi_0})$, and such that for every $x_0 \in \Ob\cE_{\pi_0}$, the functor $x_1 \longmapsto H(x_0,x_1) : \cE_1 \longrightarrow \Ens$ is $\pi$-accessible. Then we have}
\[
\cE_\pi \;=\; \cE(\pi)\ .
\]
% foot: Version 17 décembre 2025

%=====================================================================
\printedpage{160}
% head: A. Grothendieck : Dérivateurs, Ch. XVIII, Édition des Universités de Montpellier II et Paris VII

\item[\textit{3°)}] \textit{If $\pi_0$ and $\pi$ are as in 2°), and if moreover $\pi$ is useful for $\cE_1$, then $\pi$ is useful for $\cE$, hence (by virtue of 2°)}
\[
\cE \;\xleftarrow{\ \sim\ }\; \Ind_\pi(\cE(\pi))\ .
\]
\end{enumerate}
We can now state the sought-for generalization of 4.14.17:

\msp{page 422}
\textbf{Theorem 4.14.26} (generalization of 4.14.13). \textit{Let $p : \cE \longrightarrow \cC$ be a functor, with $\cC$ small. We assume that $p$ satisfies hypothesis H) of page 409, which is equivalent (using 4.14.24) to the following:}
\begin{enumerate}
\item[\textit{H$'$)}] \textit{$\forall\ s \in \Ob\cC$, the fiber category $\cE_s$ is accessible, and for every $\lambda : s \longrightarrow t$ in $\Fl\cC$, the functor}
\[
\left\{
\begin{array}{lcl}
H_\lambda : \cE_s^\circ \times \cE_t & \longrightarrow & \Ens\\[1ex]
\qquad (x_0,x_1) & \longmapsto & \Hom_\lambda(x_0,x_1)
\end{array}
\right.
\]
\textit{satisfies the following two conditions:}
\begin{enumerate}
\item[\textit{a)}] \textit{$\forall\ x_0 \in \Ob\cE_s^\circ$, $x_1 \longmapsto H_\lambda(x_0,x_1) : \cE_t \longrightarrow \Ens$ is accessible.}
\item[\textit{b)}] \textit{There exists an infinite cardinal $\pi_\lambda$, such that for every $x_1 \in \Ob\cE_t$, the functor $x_0 \longmapsto H_\lambda(x_0,x_1) : \cE_s \longrightarrow (\Ens)^\circ$ is $\pi_\lambda$-accessible.}
\end{enumerate}
\end{enumerate}
\textit{This being so, the category $\underline{\Gamma}(\cE/\cC)$ of sections of $\cE$ over $\cC$ is accessible, and the functors $\varepsilon_s : \underline{\Gamma}(\cE/\cC) \longrightarrow \cE_s$ are accessible. We can be more precise in a way similar to 4.14.18 (statement in due form left to the reader).}

\msp{page 423}
\textsc{Proof} of 4.14.26. It is carried out exactly like that of 4.14.17 (p. 400 ff.), but is conceptually simpler. We assume $\cC$ small, and introduce the categories
\[
\Pi \;=\; \prod_{s\,\in\,\Ob\cC} \cE_s\ , \qquad \Pi' \;=\; \prod_{\lambda\,\in\,\Fl\cC} \cE_\lambda\ ,
\]
where
\[
\cE_\lambda \;=\; \cE \times_{\cC} (\boldsymbol{\Delta}_1, i_\lambda)\ ,
\]
and the two canonical functors
\[
\Pi \;\underset{G}{\overset{F}{\rightrightarrows}}\; \Pi'
\]
with components
\[
\begin{aligned}
F_\lambda \;&=\; i_\lambda \circ \pr_{s(\lambda)}\\
G_\lambda \;&=\; j_\lambda \circ \pr_{b(\lambda)}\ ,
\end{aligned}
\]
where for $\lambda : s \longrightarrow t$, we set
\[
s(\lambda) \;=\; s\ , \qquad b(\lambda) \;=\; t\ ,
\]
and $i_\lambda$, $j_\lambda$ are the canonical inclusions
% In the print, the vertical arrow p_lambda and Delta_1 are enclosed together in a pair of large parentheses.
\[
\begin{tikzcd}
\cE_s \arrow[r,"i_\lambda"] \arrow[d,dash,"\wr"] & \cE_\lambda \arrow[dd,"p_\lambda"] & \cE_t\ . \arrow[l,"j_\lambda"'] \arrow[d,dash,"\wr"] \\
(\cE_\lambda)_0 & & (\cE_\lambda)_1 \\
& \boldsymbol{\Delta}_1 &
\end{tikzcd}
\]
% foot: Version 17 décembre 2025

%=====================================================================
\printedpage{161}
% head: A. Grothendieck : Dérivateurs, Ch. XVIII, Édition des Universités de Montpellier II et Paris VII

\msp{page 424}
We denote by $\cM'$ the category
\[
\cM' \;=\;
\begin{array}{l}
\text{category of pairs } (X,u)\text{, with}\\
X \in \Ob\Pi,\ u : F(x) \longrightarrow G(x),
\end{array}
\] % [sic] F(x), G(x) -> F(X), G(X), the object is named X
which is constructible by 4.14.22.1 is accessible (for $\Pi$ and $\pi'$ and the functors $F$ and $G$ are, by hypothesis H) made on $p$). We can interpret it as the category of sections % [sic] \pi' -> \Pi', the category introduced is \Pi'
of $\mathrm{Diag}(\cE)$ over $\mathrm{Diag}(\cC)$. We have a strictly fully faithful functor
\[
\cM \;\defeq\; \underline{\Gamma}(\cE/\cC) \;\hookrightarrow\; \cM'\ ,
\]
which lets us interpret $\cM$ as a strictly full subcategory of $\cM'$. We then see as in the case of 4.14.17 that this subcategory is accessible in $\cM'$ -- this is essentially sorital. It follows that $\cM$ is accessible and that the functors $\cM \longrightarrow \cE_s$ are also, as composites of
\[
\cM \;\overset{\text{\upshape acc.}}{\hookrightarrow}\; \cM' \;\overset{\text{\upshape acc.}}{\hookrightarrow}\; \Pi \;\xrightarrow{\ \ \pr_s\ \ \text{\upshape acc.}\ \ }\; \cE_s\ ,
\]

\hfill q.e.d.

\msp{page 425}

\subsection*{4.15\quad Accessible category over another}

\textbf{Proposition 4.15.1.} \textit{Let $f : \cE \longrightarrow \cC$ be a functor with $\cC$ small. We assume in addition that condition H (p. 409) is satisfied, which we make explicit here in the equivalent form H$'$:}
\begin{enumerate}
\item[\textit{H 1)}] \textit{For every $s \in \Ob\cC$, $\cE_s$ is accessible.}
\item[\textit{H 2)}] \textit{For every $\lambda \in \Fl\cC$, $\lambda : s \longrightarrow t$, there exists an infinite cardinal $\pi_\lambda$ such that for every $x_1 \in \Ob\cE_t$, the functor}
\[
x_0 \;\longmapsto\; \Hom_\lambda(x_0,x_1) \;:\; \cE_s \;\longrightarrow\; (\Ens)^\circ
\]
\textit{is $\pi_\lambda$-accessible ($\cE_s$ satisfying in addition $L_{\pi_\lambda}$).}
\item[\textit{H 3)}] \textit{For every $\lambda \in \Fl\cC$, $\lambda : s \longrightarrow t$, and for every $x_0 \in \Ob\cE_s$, the functor}
\[
x_1 \;\longmapsto\; \Hom_\lambda(x_0,x_1) \;:\; \cE_t \;\longrightarrow\; (\Ens)
\]
\textit{is accessible.}
\end{enumerate}
\textit{Then the inclusion functors $i_s : \cE_s \longrightarrow \cE$ are accessible. Moreover, if $\cC$ ``has no projectors'' (i.e. if every endomorphism $p$ of an object $s$ of $\cC$ such that $p^2 = p$, is the identity), then $\cE$ is accessible, and $f : \cE \longrightarrow \cC$ is accessible.}

This follows from the more precise result:
% foot: Version 17 décembre 2025

%=====================================================================
\printedpage{162}
% head: A. Grothendieck : Dérivateurs, Ch. XVIII, Édition des Universités de Montpellier II et Paris VII

\msp{page 426}
\textbf{Corollary 4.15.2} (\footnote{For a) and b) we use only H 1 and H 2, to the exclusion of H 3.}). \textit{The hypotheses on $f : \cE \longrightarrow \cC$ are those of 4.15.1 ($\cC$ small, and $f$ satisfies $H' =$ H 1) 2) 3)).}
\begin{enumerate}
\item[\textit{a)}] \textit{Let $s \in \Ob\cC$ and let $\pi_1$ be an infinite cardinal such that for every $\lambda : s \longrightarrow t$ (arrow of $\cC$ with source $t$), we have $\pi_1 \geqslant \pi_\lambda$ (where $\pi_\lambda$ is defined in the statement of H 2). Then the functor $i_s : \cE_s \longrightarrow \cE$ is $\pi_1$-accessible.} % [sic] source t -> source s, lambda : s -> t has source s (cf. p. 163 "de source s")
\item[\textit{b)}] \textit{We assume $\cC$ has no projectors. Let $\pi_1$ be a cardinal $\geqslant \sup(\pi_\lambda, \card\cC)$. Then $\cE$ satisfies $L_{\pi_1}$, and $f : \cE \longrightarrow \cC$ is $\pi$-accessible.}
\item[\textit{c)}] \textit{Suppose again that $\cC$ has no projectors (e.g. a ``rigid'' category, e.g. a preordered category). Let $\pi_0$, $\pi_1$ be cardinals such that:}
\begin{enumerate}
\item[\textcircled{\footnotesize $1^o$}] $\forall\ s \in \Ob\cC$, $\pi_0$ is useful for $\cE_s$ \textit{(such a $\pi_0$ exists thanks to H 1 and to the fact that $\cC$, hence $\Ob\cC$ small).}
\item[\textcircled{\footnotesize $2^o$}] $\pi_1 \geqslant \sup(\pi_0, \card\cC, \sup_{\lambda\in\Fl\cC}\pi_\lambda, \sup_{s\in\Ob\cC} \cardred \cE_{s\,\pi_0})$.
\item[\textcircled{\footnotesize $3^o$}] \textit{$\forall\lambda \in \Fl\cC$, and $x_0 \in \Ob\cE_{s,\pi_0}$, the functor $H_\lambda(x_0)$ of H 3, $x_1 \longmapsto \Hom_\lambda(x_0,x_1) : \cE_t \longrightarrow \Ens$ is $\pi_1$-accessible.}
\end{enumerate}
\textit{Let $\pi$ be a cardinal such that $\pi^{\pi_0} = \pi \geqslant \pi_1$. Then by b), $\cE$ satisfies $L_{\pi_1}$, hence $L_\pi$, and $f : \cE \longrightarrow \cC$ is $\pi_1$-accessible. Moreover,}
\[
i_s(\cE_{s\,\pi}) \;\subset\; \cE_\pi\ ,
\]
\textit{or again}
\[
\cE(\pi) \;\subset\; \cE_\pi\ , \tag*{$(*)$}
\]
\textit{where $\cE(\pi)$ is the full subcategory of $\cE$ such that}
\[
\Ob\cE(\pi) \;=\; \bigcup_{s\in\Ob\cC} \Ob\,(\cE_{s\,\pi})\ .
\]
\textit{As $\cE_\pi$ satisfies $L_{\pi_1}$ by b), hence a fortiori $L_\pi$, we therefore have by $(*)$ a functor} fully faithful
\[
\Ind_\pi(\cE_{s\,\pi}) \;\xrightarrow{\ \sim\ }\; \cE\ ,
\]
\textit{and this functor is an equivalence of categories (hence $\cE$ is accessible).}
\item[\textit{d)}] \textit{In fact, under the hypotheses of c) on $\cC$, $\pi_0$, $\pi_1$, $\pi$, we have $\cE(\pi) = \cE_\pi$ (not only $\cE(\pi) \subset \cE_\pi$).}
\end{enumerate}

\msp{page 427}
\textsc{Proof of 4.15.1 and 4.15.2.} It is clear that 4.15.1 is contained in 4.15.2 a) and c). We will therefore prove 4.15.2.

a) Let $(x_i)_{i\in I}$ be a direct system in $\cE_s$, with $I$ large relative to $\pi_0$, and $x$ its direct limit in $\cE_s$ (which exists, for $\cE_s$ satisfies $\mathrm{L}_{\pi_\lambda}$, hence $\mathrm{L}_{\pi_0}$). We must prove that it is a limit in $\cE$, hence that for every $y \in \Ob\cE$, the canonical arrow
\[
\Hom_{\cE}(x,y) \;\longrightarrow\; \varprojlim_{I} \Hom_{\cE}(x_i,y) \tag*{$(*)$}
\]
% foot: Version 17 décembre 2025

%=====================================================================
\printedpage{163}
% head: A. Grothendieck : Dérivateurs, Ch. XVIII, Édition des Universités de Montpellier II et Paris VII
is an isomorphism. Now $\Hom_{\cE}(x,y)$ and the $\Hom_{\cE}(x_i,y)$ are objects of $\Ens/\Hom_{\cC}(s,t)$, where $s = f(x)$, $t = f(y)$, and the transition arrows are arrows in $\Ens/\Hom_{\cC}(s,t)$. This shows that the bijectivity of $(*)$ comes down to that of
\[
\Hom_\lambda(x,y) \;\longrightarrow\; \varprojlim_{I} \Hom_\lambda(x_i,y) \tag*{$(*)$}
\]
for every $\lambda \in \Hom_{\cC}(x,y)$ \texttt{[rather} $\Hom_{\cC}(f(x),f(y))$\texttt{]}. For $\lambda$ fixed, and $y$ fixed in $\Ob\cE_t$, the bijectivity in $(*)$, for every direct system large relative to $\pi_1$ in $\cE_s$, with limit $x$, means that the functor
\[
x \;\longmapsto\; \Hom_\lambda(x,y) \;:\; \cE_s \;\longrightarrow\; (\Ens)^\circ
\]
is $\pi_0$-accessible. Now this is indeed the case, by the hypothesis on $\pi_1$, whatever the arrow $\lambda : s \longrightarrow t$ with source $s$, and whatever $y \in \Ob\cE_t$.

\msp{page 428}
Let $(x_i)_{i\in I}$ be a direct system large relative to $\pi_1$, and $(s_i)_{i\in I}$ the image of this system in $\cC$. We will prove below the

% In the left margin, beside the statement of Lemme 4.15.3 and beside the following paragraph "Mais attention ...", the edition prints a wavy marginal sign (one for each).
\textbf{Lemma 4.15.3} (\footnote{Cf. proof page 452.},\footnote{subject to the validity of 4.8.13 (But this is now established, in a sufficient form, in 4.15.10).}). \textit{Let $\cC$ be a category, $\pi$ a cardinal, such that} $\card\cC \leqslant \pi$\textit{, and that $\cC$ has no projectors. Then for every direct system $(s_i)_{i\in I}$ large relative to $\pi$ in $\cC$, there exists a cofinal subset $J$ in $I$ such that the induced system $(s_i)_{i\in J}$ is} constant, \textit{i.e. such that the $s_i$ are equal to an $s \in \cC$, and that the transition morphisms are all equal to} $\id_s$.

But watch out, we rely in proving 4.15.3 on theorem 4.8.13 (page 212-216), which I still have not proved. In the case where $\cC$ is rigid (i.e. $\End_{\cC}(s) = \{\id_s\}\ \forall\ s \in \Ob\cC$), the proof is however immediate, cf. below.

\msp{page 429}
To prove that $x = \varinjlim\limits_{I} x_i$ exists in $\cE$, and commutes with $f$, we can replace $I$ by a cofinal subcategory. This brings us back, thanks to the lemma, to the case where $(s_i)_{i\in I}$ is constant, hence to the case where the direct system comes from a direct system in a fiber category $\cE_s$. But by a), we know that its limit $x$ in $\cE_s$ is a limit in $\cE$, whence the existence of the latter, and commutation with $f$, since $\varinjlim s_i = s$ ( $\varinjlim$ of a constant filtering direct system).

c) We must prove that if $x \in \Ob\cE_s$ is $\pi$-accessible in $\cE_s$, it is so in $\cE$, i.e. that for every direct system $(y_i)_{i\in I}$ large relative to $\pi$, the canonical arrow
\[
\varinjlim_{i} \Hom(x,y_i) \;\longrightarrow\; \Hom(x,y) \tag*{$(**)$}
\]
is bijective. Now as by 2°) $\pi \geqslant \pi_1 \geqslant \sup(\card S, \sup_\lambda \pi_\lambda)$, it follows again from the lemma % [sic] card S -> card \mathcal{C}, no S is defined; the condition 2° bounds card \mathcal{C}
that, possibly replacing $I$ by a cofinal subcategory, we can assume that the direct system $(y_i)_{i\in I}$ comes from a system in a fiber category $\cE_t$, and that $y$ is its $\varinjlim$ in $\cE_t$. But then the sets $\Hom(x,y)$ and $\Hom(x,y_i)$, and the transition morphisms between them, are again in $\Ens/\Hom(s,t)$, and we are reduced
% foot: Version 17 décembre 2025

\end{document}
